\expandafter\def\csname HistoricalNoteText4\endcsname{Diese ganze Sache wird durch eine andere, bereits im Druck befindliche Untersuchung sehr erhellt werden. Dort hätte ich bei einem zwar ganz anderen, aber dennoch analogen Gegenstand eine ähnliche Freiheit mit genau demselben Recht gebrauchen können, wie es hier bei den Gleichungen von allen Analysten geschehen ist. Obwohl ich mithilfe solcher Fiktionen die Beweise mehrerer Wahrheiten mit wenigen Worten hätte abschließen können, die ohne sie sehr schwierig wurden und die feinsten Kunstgriffe erfordern, habe ich doch vorgezogen, ganz darauf zu verzichten, und hoffe, dass ich nur wenige zufriedengestellt hätte, wenn ich die Methode der Analysten nachgeahmt hätte.}

\expandafter\def\csname HistoricalNoteText5\endcsname{Euler hat irrtümlich $C$ und setzt daher auch später fälschlich $pqr=C$.}

\expandafter\def\csname HistoricalNoteText6\endcsname{In diese Darstellung scheint sich ein Fehler eingeschlichen zu haben: Auf S. 118, Zeile 5, muss man nämlich anstelle des Zeichens $p$ (``man wählte nur diejenigen, in denen $p$ vorkam'', usw.) notwendig \textit{ein und dieselbe beliebige Wurzel der vorgelegten Gleichung} oder etwas Ähnliches lesen, da das Erstere keinen Sinn hat.}
